\documentclass[11pt]{article}
\usepackage[a4paper,margin=25mm]{geometry}
\usepackage{amsmath,amssymb,amsthm,booktabs,microtype}
\usepackage[T1]{fontenc}
\usepackage{lmodern}
\usepackage[dvipsnames]{xcolor}
\usepackage[colorlinks=true,linkcolor=MidnightBlue,citecolor=MidnightBlue,urlcolor=MidnightBlue]{hyperref}
\usepackage{fancyhdr}
\pagestyle{fancy}\fancyhf{}
\fancyhead[L]{\small The upper saturation spectrum of Berge triangles}
\fancyhead[R]{\small Dev D. Goyal}
\fancyfoot[C]{\thepage}
\setlength{\headheight}{14pt}
\hypersetup{pdftitle={The upper saturation spectrum of Berge triangles in triple systems},
pdfauthor={Dev D. Goyal},
pdfsubject={Fully autonomous agentic research},
pdfkeywords={Berge triangle, saturation spectrum, triple systems, 3-uniform hypergraphs, extremal combinatorics, stability, Lean 4}}
\newtheorem{theorem}{Theorem}
\newtheorem{lemma}[theorem]{Lemma}
\newtheorem{corollary}[theorem]{Corollary}
\theoremstyle{remark}
\newtheorem{remark}[theorem]{Remark}
\newcommand{\F}{F}
\newcommand{\sh}{\partial}
\title{\textbf{The upper saturation spectrum of Berge triangles\\in triple systems}}
\author{Dev D. Goyal}
\date{31 October 2026}
\begin{document}
\maketitle
\begin{abstract}
Let $E_n=\lfloor n^2/8\rfloor$ be Gy\H{o}ri's maximum number of edges in a
Berge-triangle-free triple system on $n$ vertices, and let
$L_n=\lfloor (n-1)^2/8\rfloor$. We show that for every $n\ge 8$ every
Berge-triangle-free triple system with more than $L_n$ edges is a spanning
subsystem of a paired product $\F(n,p)$: $p$ disjoint pairs, $n-2p$ further
vertices, and all triples formed by one pair and one further vertex.
Consequently every saturated system with more than $L_n$ edges is a paired
product, and the part of the Berge-triangle saturation spectrum above $L_n$
is exactly $\{p(n-2p):1\le p<n/2\}\cap(L_n,E_n]$. The threshold is sharp:
for every $n\ge 8$ there is a saturated system with exactly $L_n$ edges that
is not a paired product. The upper window has $\lfloor n/4\rfloor$ or
$\lfloor n/4\rfloor+1$ integers and contains $n/4-O(\sqrt n)$ sizes that are
not in the spectrum, given explicitly by quadratic deficits. The proof uses
the own-pair reduction of Frankl, F\"uredi, Goorevitch, Holzman and
Simonyi, a rigidity lemma for dense bipartite own-pair graphs, and the
equality case of Erd\H{o}s's bound for nonbipartite triangle-free graphs.
A Lean~4 project proves the definitions-level facts that every paired
product is Berge-triangle-free and saturated with $p(n-2p)$ edges; the
containment theorem has a conventional proof, and exact checkers verify the
constructions and the arithmetic for small orders.
\end{abstract}

\section{Definitions and results}

A \emph{triple system} on a vertex set $V$ is a set $H$ of $3$-element
subsets of $V$, its \emph{edges}. All vertices of $V$ count, including
isolated ones. $H$ contains a \emph{Berge triangle} if there are distinct
vertices $x,y,z$ and distinct edges $h_1,h_2,h_3\in H$ with
$\{x,y\}\subseteq h_1$, $\{y,z\}\subseteq h_2$ and $\{z,x\}\subseteq h_3$;
copies need not be induced. $H$ is \emph{saturated} if it contains no
Berge triangle but $H\cup\{T\}$ contains one for every $3$-set $T\notin H$.
The \emph{saturation spectrum} $\operatorname{ES}_3(n,\text{Berge-}K_3)$ is
the set of sizes $|H|$ of saturated systems on $n$ vertices.

Gy\H{o}ri~\cite{Gyori} proved that a Berge-triangle-free triple system on
$n$ vertices has at most $E_n=\lfloor n^2/8\rfloor$ edges, and Frankl,
F\"uredi, Goorevitch, Holzman and Simonyi~\cite{FFGHS} characterized the
extremal systems. Saturation numbers of Berge triangles were studied
in~\cite{EGGMS}. Bushaw, English, Heath, Johnston and
Rombach~\cite[Section~6]{BEHJR}, after determining spectra of Berge stars,
proposed the spectrum of the Berge triangle as a further problem, recording
its minimum $\lceil (n-1)/2\rceil$ and maximum $E_n$, and suspected that
only finitely many values near the top are missing. Their
remark is a suspicion rather than a numbered conjecture.

For $1\le p<n/2$ the \emph{paired product} $\F(n,p)$ has vertex set
$P_1\cup\dots\cup P_p\cup B$, where the $P_i$ are disjoint pairs and
$|B|=n-2p\ge1$, and edge set $\{P_i\cup\{b\}:1\le i\le p,\ b\in B\}$. It has
$p(n-2p)$ edges; for suitable $p$ this attains $E_n$. Write
\[
E_n=\Big\lfloor\frac{n^2}{8}\Big\rfloor,\qquad
L_n=\Big\lfloor\frac{(n-1)^2}{8}\Big\rfloor,\qquad W_n=E_n-L_n .
\]

\begin{theorem}[dense containment]\label{thm:containment}
Let $n\ge8$. Every Berge-triangle-free triple system $H$ on $n$ vertices with
$|H|>L_n$ is a spanning subsystem of a paired product $\F(n,p)$.
\end{theorem}

\begin{theorem}[upper spectrum]\label{thm:spectrum}
Let $n\ge 8$. A saturated triple system on $n$ vertices with more than $L_n$
edges is isomorphic to a paired product. Hence
\[
\operatorname{ES}_3(n,\text{Berge-}K_3)\cap(L_n,E_n]
=\{p(n-2p):1\le p<n/2\}\cap(L_n,E_n].
\]
\end{theorem}

\begin{theorem}[sharpness]\label{thm:sharp}
For every $n\ge8$ the largest saturated triple system on $n$ vertices that
is not isomorphic to a paired product has exactly $L_n$ edges. In
particular the strict threshold in Theorems~\ref{thm:containment}
and~\ref{thm:spectrum} cannot be lowered.
\end{theorem}

\begin{corollary}[explicit deficits]\label{cor:deficits}
Let $n\ge8$ and $0\le d<W_n$. Then $E_n-d$ belongs to the saturation
spectrum if and only if
\begin{itemize}
\item $d=2j^2$ for some integer $j\ge0$, when $n\equiv0\pmod 4$;
\item $d=2j(j+1)$ for some integer $j\ge0$, when $n\equiv2\pmod 4$;
\item $d=j(j+1)/2$ for some integer $j\ge0$, when $n$ is odd.
\end{itemize}
Here $W_n=\lfloor n/4\rfloor$, except that $W_n=\lfloor n/4\rfloor+1$ when
$n\equiv3\pmod4$. The window $(L_n,E_n]$ therefore contains
$n/4-O(\sqrt n)$ sizes outside the spectrum.
\end{corollary}

For example, above $L_{16}=28$ only $30$ and $32$ occur, and for $n=400$
only $8$ of the $100$ sizes $19901,\dots,20000$ occur
(Table~\ref{tab:examples}). The number of missing sizes is thus unbounded,
and they are not confined to a bounded distance from the maximum. This
rules out the natural uniform-in-$n$ reading of the suspicion
in~\cite{BEHJR}. The spectrum at or below $L_n$ is not determined here.

\paragraph{Prior work.} The paired products, their size, the bound $E_n$
and the classification of the extremal systems are
known~\cite{Gyori,FFGHS}. Lemma~\ref{lem:own} below is the own-pair
reduction of~\cite[Observation~2]{FFGHS}. The nonbipartite bound in
Lemma~\ref{lem:erdos} is due to Erd\H{o}s~\cite[Lemma~1]{Erdos}; we include
its proof because we need its equality case. Gerbner~\cite{Gerbner}
obtained qualitative matching and shadow structure for near-extremal
Berge-triangle-free systems; such structure may already imply,
after a saturation argument, that the spectrum misses many sizes near the
top. What we add is the exact finite threshold $L_n$ for every $n\ge8$,
the nonproduct construction at the threshold, and the explicit spectrum in
the whole window. We did not find these statements in the searched
literature. The full text of a 2026 preprint of Xu, Zeng and Zhang on the
Tur\'an number of a matching together with a Berge triangle
(arXiv:2607.12300) could not be inspected, and its structural lemmas may
bear on Theorem~\ref{thm:containment}. No priority claim is made.

\section{Own pairs}

A pair of vertices is an \emph{own pair} of an edge $h$ if $h$ is the only
edge containing it. The \emph{shadow} $\sh H$ is the graph of all pairs
contained in some edge.

\begin{lemma}[{\cite[Observation~2]{FFGHS}}]\label{lem:own}
Let $H$ be Berge-triangle-free. Every edge has at least two own pairs.
Choosing exactly two own pairs from each edge gives a triangle-free graph $G$
with $e(G)=2|H|$. Every triangle of $\sh H$ is an edge of $H$.
\end{lemma}

\begin{proof}
If $h=\{a,b,c\}$ had non-own pairs $ab$ and $ac$, pick edges
$h_1\ne h\ni a,b$ and $h_2\ne h\ni a,c$. They are distinct, since an edge
containing $ab$ and $ac$ equals $h$, and $h_1,h_2,h$ with the pairs
$ab,ac,bc$ form a Berge triangle. Own pairs chosen from different edges are
distinct, so $e(G)=2|H|$. If $abc$ is a triangle of $\sh H$ but not an edge,
the edges containing $ab$, $bc$, $ca$ are pairwise distinct (a common edge
of two of them would be $abc$) and form a Berge triangle. Finally, a
triangle of $G$ is a triangle of $\sh H$, hence an edge $h$; but only two of
its pairs were chosen from $h$, and the third, being own, cannot be chosen
from any other edge.
\end{proof}

Fix such a \emph{selected own-pair graph} $G$, and let $M$ be the set of
\emph{omitted pairs}, the third pair of each edge. A selected pair is own
and is selected only from its own edge, where it is not omitted, so
$M\cap G=\emptyset$ and $\sh H=G\cup M$. Omitted pairs may have codegree one.

\section{Dense bipartite own-pair graphs}

\begin{lemma}[bipartite rigidity]\label{lem:rigid}
Let $H\neq\emptyset$ be Berge-triangle-free with a selected own-pair graph
$G$ that is bipartite with parts $A,B$, $a=|A|$, $b=|B|$, and suppose
$D=ab-e(G)<\min(a,b)$. Then $H$ is a spanning subsystem of a paired product.
\end{lemma}

\begin{proof}
The two selected pairs of an edge form a path of length two in $G$, so each
omitted pair lies inside $A$ or inside $B$.

\emph{$M$ is a matching.} Suppose $xy,xz\in M$ lie in $A$. If some $w\in B$
were $G$-adjacent to $x,y,z$, then $xyw$ and $xzw$ would be triangles of
$\sh H$, hence edges by Lemma~\ref{lem:own}, and the selected pair $xw$
would lie in two edges. So each $w\in B$ misses one of $xw,yw,zw$, and
$D\ge b$, a contradiction. The same holds inside $B$.

Let $p$ and $q$ be the numbers of pairs of $M$ inside $A$ and $B$, and
$r=a-2p$, $s=b-2q$ the numbers of unmatched vertices.

\emph{Blocks.} Let $P\subseteq A$ and $Q\subseteq B$ be pairs of $M$. An edge
supplying a $G$-edge between $P$ and $Q$ has its omitted pair in $M$ meeting
$P$ or $Q$, so it is $P\cup\{v\}$ with $v\in Q$ or $Q\cup\{u\}$ with $u\in P$,
and contributes two $G$-edges to this $2\times2$ block. Two such edges of
opposite kinds share a cross pair, which is selected and own; two of the same
kind, say $P\cup\{u\}$ and $P\cup\{v\}$ with $Q=\{u,v\}$, together with an
edge $h\supseteq Q$ (which exists since $Q\in M$) and any $x\in P$, give the
Berge triangle on $x,u,v$. Hence each block has at most two $G$-edges.
An unmatched $u\in A$ and an unmatched $v\in B$ are not $G$-adjacent: the
edge supplying $uv$ would have an omitted pair at $u$ or $v$. Counting
these disjoint sets of missing cross pairs,
\[
D\ge 2pq+rs .
\]
If $p,q\ge1$ and $r\le s$ then $rs\ge r$ and $2pq\ge 2p$, so $D\ge a$;
symmetrically $D\ge b$ if $s\le r$. So $p$ or $q$ is zero, and since
$M\ne\emptyset$ we may assume $p>0=q$. Then every omitted pair lies in $A$,
and an unmatched vertex of $A$ would have $G$-degree zero, forcing $D\ge b$.
So $r=0$: $A$ is partitioned by the pairs of $M$, and every edge is a pair of
$M$ together with a vertex of $B$. As $b>D\ge0$, this is a paired product
containing $H$.
\end{proof}

\section{The nonbipartite threshold and its equality case}

Let $T_n=\lfloor (n-1)^2/4\rfloor+1$.

\begin{lemma}[{Erd\H{o}s~\cite[Lemma~1]{Erdos}}]\label{lem:erdos}
A nonbipartite triangle-free graph on $n$ vertices has at most $T_n$ edges.
\end{lemma}

\begin{proof}
Let $C$ be a shortest odd cycle, of length $2k+1\ge5$; it has no chords. A
vertex outside $C$ has at most two neighbours on $C$: three neighbours split
$C$ into gaps of length at least two, one gap is odd and at most $2k-3$, and
it closes a shorter odd cycle through the vertex. With $m=n-2k-1$ and
Mantel's theorem outside $C$,
\[
e(G)\le (2k+1)+2m+\Big\lfloor\frac{m^2}{4}\Big\rfloor
=T_n-(k-2)(n-k-1)\le T_n,
\]
where the identity uses $\lfloor (m+2k)^2/4\rfloor=\lfloor m^2/4\rfloor+km+k^2$.
\end{proof}

\begin{lemma}\label{lem:equality}
Let $n=4t+3\ge11$, and let $G$ be nonbipartite and triangle-free with
$e(G)=T_n$. Then $V(G)$ is partitioned into $A,D,C,\{u\},\{v\}$ with
$|C|=2t$, $|A|+|D|=2t+1$, $|A|,|D|\ge1$, and the edges of $G$ are exactly
those between $C$ and $A\cup D$, those between $u$ and $A$, those between
$v$ and $D$, and $uv$.
\end{lemma}

\begin{proof}
Equality forces $k=2$, two neighbours on the $5$-cycle $c_0\dots c_4$ for
every outside vertex, and equality in Mantel's theorem outside the cycle;
so the $4t-2$ outside vertices span a complete bipartite graph with parts
$U,V$ of size $2t-1$. An outside vertex adjacent to $c_{i-1},c_{i+1}$ has
\emph{type} $i$. For adjacent $x\in U$, $y\in V$ the neighbourhoods on the
cycle are disjoint, which happens exactly when their types are adjacent
indices of the cycle. As the $5$-cycle contains no $K_{2,2}$, one side, say
$U$, has a single type, which we rotate to $0$. Then $V$ splits into
$V_1,V_4$ by type, and $C=\{c_0\}\cup U$, $A=\{c_1\}\cup V_1$,
$D=\{c_4\}\cup V_4$, $u=c_2$, $v=c_3$ have the stated properties.
\end{proof}

\begin{lemma}[switching]\label{lem:switch}
If a Berge-triangle-free $H$ has a selected own-pair graph of the form in
Lemma~\ref{lem:equality}, then another selection of own pairs gives a
bipartite graph with the same number of edges.
\end{lemma}

\begin{proof}
Let $h$ be the edge supplying $uv$. Its other selected pair contains $u$ or
$v$; by symmetry $h=\{u,v,x\}$ with $ux$ selected, so $x\in A$ and $vx$ is
omitted. The pair $vx$ is also own: another edge containing $vx$ would have
$vx$ as its omitted pair and a common $G$-neighbour of $v$ and $x$ as third
vertex, and $u$ is their only common neighbour. Replacing $uv$ by $vx$
keeps two own pairs per edge and gives a graph bipartite with parts
$A\cup D$ and $C\cup\{u,v\}$.
\end{proof}

\section{Proofs of Theorems~\ref{thm:containment} and~\ref{thm:spectrum}}

\begin{proof}[Proof of Theorem~\ref{thm:containment}]
Let $|H|>L_n$ and select an own-pair graph $G$; it is triangle-free with
$e(G)=2|H|$. Comparing residues of $(n-1)^2$ modulo $8$ gives $2|H|\ge
2L_n+2\ge T_n$, with equality only if $n\equiv3\pmod4$ and $|H|=L_n+1$. If
$e(G)>T_n$, then $G$ is bipartite by Lemma~\ref{lem:erdos}; if $e(G)=T_n$
and $G$ is not bipartite, Lemmas~\ref{lem:equality} and~\ref{lem:switch}
give a bipartite selection with the same number of edges. So some selected
own-pair graph $G$ is bipartite with $e(G)\ge T_n$. If its parts have sizes
$a\le b$, then $a(b-1)\le\lfloor (n-1)^2/4\rfloor$, so
$D=ab-e(G)\le ab-T_n\le a-1<\min(a,b)$, and Lemma~\ref{lem:rigid} applies.
\end{proof}

\begin{lemma}\label{lem:product}
Every paired product $\F(n,p)$ is Berge-triangle-free and saturated.
\end{lemma}

\begin{proof}
Each edge has exactly one vertex in $B$. A Berge triangle cannot contain two
vertices of $B$, since one edge would contain both. If it contains none,
the pair $xy$ forces $x,y$ into one pair $P_i$, and likewise for $yz$, so
three distinct vertices lie in $P_i$. If it contains one, say $z\in B$, then
$x,y\in P_i$ and both edges through $z$ equal $P_i\cup\{z\}$. For saturation,
an absent triple contains two vertices of $B$, which are joined by the
two-edge Berge path $b,x,b'$ through any $x\in P_1$, or vertices from two
different pairs, joined through any $b\in B$; the only other triples,
one pair and one vertex of $B$, are present. Adding the triple closes the
path into a Berge triangle.
\end{proof}

\begin{proof}[Proof of Theorem~\ref{thm:spectrum}]
If $H$ is saturated with $|H|>L_n$, then $H\subseteq \F$ for a paired product
$\F$ by Theorem~\ref{thm:containment}. Adding any triple of $\F\setminus H$
to $H$ stays inside the Berge-triangle-free $\F$, so saturation forces
$H=\F$. Conversely every $\F(n,p)$ is saturated by Lemma~\ref{lem:product}.
\end{proof}

\section{The sharpness construction}

For $p\ge1$ and $q\ge3$ take disjoint pairs $P_1=\{x,y\},P_2,\dots,P_p$, a
set $S$ of $q-2$ vertices and two further vertices $u,v$, so $n=2p+q$. Let
$J(p,q)$ consist of
\[
P_i\cup\{z\}\ (1\le i\le p,\ z\in S),\qquad
P_i\cup\{v\}\ (2\le i\le p),\qquad \{x,u,v\}.
\]
It has $p(q-2)+(p-1)+1=p(q-1)$ edges.

\begin{lemma}\label{lem:J}
$J(p,q)$ is saturated. For $p\ge2$ it is not isomorphic to a paired product.
\end{lemma}

\begin{proof}
\emph{Freeness.} Without $\{x,u,v\}$ the system is contained in the paired
product with pairs $P_i$ and base $S\cup\{u,v\}$, so a Berge triangle must
use $\{x,u,v\}$ for a pair inside it, and the other two edges form a
two-edge Berge path between the two vertices of that pair. Since $u$ lies
in no other edge, the pair is $xv$. A first edge at $x$ is $P_1\cup\{z\}$
with $z\in S$, and neither $y$ nor $z$ lies in an edge with $v$.

\emph{Saturation.} The vertex $u$ has a two-edge Berge path to every vertex
outside $\{x,u,v\}$: to $y$ or $z\in S$ through $x$, using $\{x,u,v\}$ and
$P_1\cup\{z\}$, and to a vertex of $P_i$, $i\ge2$, through $v$, using
$\{x,u,v\}$ and $P_i\cup\{v\}$. An absent triple containing $u$ contains
such a vertex. For an absent triple avoiding $u$: two vertices in different
pairs are joined through any $z\in S$; two vertices of $S$ are joined
through $x$, and so are $v$ and $z\in S$, using $\{x,u,v\}$ and
$P_1\cup\{z\}$. Otherwise the triple is a pair $P_i$ with one vertex of
$S\cup\{v\}$; all of these are present except $P_1\cup\{v\}$, whose pair
$yv$ closes the path $y,x,v$.

\emph{Not a product.} The vertex $u$ has degree one, and for $p\ge2$ the
edges $P_1\cup\{z\}$ and $P_2\cup\{v\}$ are disjoint. In a paired product a
vertex of degree one forces $p=1$ or $|B|=1$, and then all edges share a
vertex.
\end{proof}

\begin{proof}[Proof of Theorem~\ref{thm:sharp}]
Take $p=\lfloor n/4\rfloor$ and $q=n-2p$. For $n\ge8$ we have $p\ge2$ and
$q\ge3$, and checking the four residues of $n$ modulo $4$ gives
$p(q-1)=p(n-2p-1)=L_n$. So $J(p,q)$ is a saturated system with $L_n$ edges
that is not a paired product. By Theorem~\ref{thm:spectrum} no larger
nonproduct saturated system exists. Since $J(p,q)$ is saturated, it is not
contained in any paired product (saturation would force equality), so the
threshold of Theorem~\ref{thm:containment} is sharp as well.
\end{proof}

The construction proves the value $L_n$ outside the product class; it does
not classify the nonproduct saturated systems with $L_n$ edges, and other
such systems exist, for instance at $n=11$.

\section{Arithmetic of the window}

Write $n=4t+r$. Parametrizing $p$ around $n/4$ gives
\begin{align*}
n=4t,\ p=t+j:&\quad E_n-p(n-2p)=2j^2,\\
n=4t+2,\ p=t+j:&\quad E_n-p(n-2p)=2j(j-1),\\
n=4t+1,\ p=t+j:&\quad E_n-p(n-2p)=2j^2-j,\\
n=4t+3,\ p=t+1+j:&\quad E_n-p(n-2p)=2j^2+j,
\end{align*}
with $j\in\mathbb Z$. As $j$ ranges over the integers, the nonnegative values
of $2j(j-1)$ are the numbers $2k(k+1)$, and those of $2j^2\mp j$ are the
triangular numbers $k(k+1)/2$, $k\ge0$. If $d<W_n$ then $E_n-d>L_n>0$, so a
solution of $p(n-2p)=E_n-d$ automatically has $0<p<n/2$. With
Theorem~\ref{thm:spectrum} this proves Corollary~\ref{cor:deficits}. The
values $E_n$, $L_n$ are $2t^2$, $2t^2-t$ for $n=4t$; $2t^2+t$, $2t^2$ for
$n=4t+1$; $2t^2+2t$, $2t^2+t$ for $n=4t+2$; and $2t^2+3t+1$, $2t^2+2t$ for
$n=4t+3$, which gives $W_n$. Each quadratic family meets the window in
$O(\sqrt{W_n})$ values.

\begin{table}[ht]
\centering
\begin{tabular}{rrrll}
\toprule
$n$ & $L_n$ & $E_n$ & sizes in $(L_n,E_n]$ in the spectrum & missing sizes\\
\midrule
8 & 6 & 8 & 8 & 7\\
11 & 12 & 15 & 14, 15 & 13\\
16 & 28 & 32 & 30, 32 & 29, 31\\
40 & 190 & 200 & 192, 198, 200 & 191, 193--197, 199\\
400 & 19900 & 20000 & 19902, 19928, 19950, 19968, & 92 of the 100 sizes\\
 & & & 19982, 19992, 19998, 20000 & \\
\bottomrule
\end{tabular}
\caption{The upper window for some orders.}\label{tab:examples}
\end{table}

\section{Verification}\label{sec:verification}

\paragraph{Exact checks.} The program \texttt{check\_spectrum.py} in
\texttt{certificate/}, which uses only the Python standard library, builds
$J(\lfloor n/4\rfloor, n-2\lfloor n/4\rfloor)$ and every paired product above the threshold for
$8\le n\le 20$, and checks freeness, saturation, sizes and product
recognition by direct search. It compares the product sizes in
$(L_n,E_n]$ with Corollary~\ref{cor:deficits} and the formula for $W_n$
for $8\le n\le 2000$, and recomputes Table~\ref{tab:examples}. It also
checks four witness systems directly: $J(2,7)$ at $n=11$, $J(4,8)$ at
$n=16$, and paired products with $30$ and $32$ edges at $n=16$. A separate checker validates the stored
search-free certificates of these witnesses, listing for every absent
triple a closing Berge path.

\paragraph{Lean.} The Lean~4.35.0 project in \texttt{lean/}, pinned to
Mathlib commit \texttt{56a0058}, defines triple systems on an arbitrary
vertex type, Berge triangles, freeness, saturation and paired products as
in Section~1. It proves that every paired product is Berge-triangle-free
and saturated, that its number of edges is the number of pairs times the
size of the base, and that on a finite vertex type twice the number of
pairs plus the size of the base is the number of vertices. The audited
declarations depend only on \texttt{propext}, \texttt{Classical.choice}
and \texttt{Quot.sound}. Theorem~\ref{thm:containment}, the sharpness
construction and the arithmetic of the window have conventional proofs
only in this version.

\section{Limitations and open questions}

The spectrum at or below $L_n$ is open, and the nonproduct saturated
systems with exactly $L_n$ edges are not classified. For $n\le7$ the
statements are not claimed. The literature comparison is bounded, as
described in Section~1.

\section*{Research process and AI disclosure}

This is fully autonomous agentic research. Dev D. Goyal is responsible for
publication. No expert validation or external peer review is claimed, and
the novelty assessment is provisional.

The problem selection, the theorems, their proofs, the construction, the
exact checkers, the witness certificates and the literature search were
produced by autonomous AI research agents in an earlier session, which also
ran internal AI-assisted reviews of the proof; these reviews are not peer
review. Its records do not name the model used for the research; they
describe the internal reviews as Astra reviews. During that session the
target threshold was strengthened from $E_n-\lfloor n/4\rfloor$ to $L_n$,
which differs from it by one edge when $n\equiv3\pmod 4$, and the
containment statement was extended from saturated to all free systems.

Opus 5.5 by Anthropic used with Claude Code assessed the package for this
release. It checked each proof step, wrote its own exact checker for the
constructions and the arithmetic, reran the certificate checker, wrote the
Lean project and ran its fresh build and axiom audit, and wrote this
manuscript. AI checks are not human peer review.

\begin{thebibliography}{9}
\bibitem{BEHJR}
N.~Bushaw, S.~English, E.~Heath, D.~Johnston and P.~Rombach,
\emph{The saturation spectrum of Berge stars}, preprint (2025),
\href{https://arxiv.org/abs/2502.17686}{arXiv:2502.17686}.

\bibitem{EGGMS}
S.~English, D.~Gordon, N.~Graber, A.~Methuku and E.~C.~Sullivan,
\emph{Saturation of Berge hypergraphs}, Discrete Mathematics \textbf{342}
(2019), 1738--1761.
\href{https://doi.org/10.1016/j.disc.2019.01.031}{doi:10.1016/j.disc.2019.01.031}.
\href{https://arxiv.org/abs/1710.03735}{arXiv:1710.03735}.

\bibitem{Erdos}
P.~Erd\H{o}s, \emph{On a theorem of Rademacher--Tur\'an}, Illinois Journal
of Mathematics \textbf{6} (1962), 122--127.

\bibitem{FFGHS}
P.~Frankl, Z.~F\"uredi, A.~Goorevitch, R.~Holzman and G.~Simonyi,
\emph{Triangle-free triple systems}, Electronic Journal of Combinatorics
\textbf{33}(2) (2026), \#P2.34.
\href{https://doi.org/10.37236/14115}{doi:10.37236/14115}.
Preprint \href{https://arxiv.org/abs/2405.16452}{arXiv:2405.16452}.

\bibitem{Gerbner}
D.~Gerbner, \emph{The Tur\'an number of Berge book hypergraphs}, SIAM
Journal on Discrete Mathematics \textbf{38} (2024), 2896--2912.
\href{https://arxiv.org/abs/2111.11162}{arXiv:2111.11162}.

\bibitem{Gyori}
E.~Gy\H{o}ri, \emph{Triangle-free hypergraphs}, Combinatorics,
Probability and Computing \textbf{15} (2006), 185--191.
\end{thebibliography}
\end{document}
